\section{Implementación y Pruebas}

% ============================================================
\subsection{Arquitectura de Despliegue del Pipeline Analítico}
% ============================================================

El pipeline analítico de Bouquet se despliega de manera independiente sobre un
servidor privado virtual (VPS), separado de la infraestructura transaccional
alojada en Supabase y de la capa de presentación desplegada en Vercel. Esta
separación responde a la necesidad de aislar las cargas de procesamiento batch
de las operaciones OLTP en tiempo real, evitando que las tareas de extracción,
transformación y agregación de datos compitan por recursos con las consultas
operativas del restaurante.

\subsubsection{Especificaciones del Servidor}

El entorno de ejecución reside en un VPS con las siguientes características
técnicas:

\begin{itemize}[leftmargin=1.5em]
  \item \textbf{Sistema operativo}: Ubuntu 25.10 (Plucky Puffin), kernel Linux
  actualizado.
  \item \textbf{CPU}: 2 núcleos virtuales AMD EPYC, arquitectura x86\_64, con
  virtualización KVM.
  \item \textbf{Memoria RAM}: 8 GB, suficiente para la ejecución concurrente del
  orquestador FastAPI y los trabajos de Spark en modo local.
  \item \textbf{Almacenamiento}: SSD con partición dedicada en \texttt{/data} para
  los archivos Parquet generados por el pipeline.
  \item \textbf{Red}: interfaz pública con firewall administrado mediante
  \texttt{iptables} y \texttt{fail2ban}, más una red interna Docker para
  comunicación entre contenedores sin exponer puertos de base de datos.
\end{itemize}

\subsubsection{Contenerización con Docker}

La totalidad del pipeline se ejecuta dentro de contenedores Docker, lo que
garantiza reproducibilidad del entorno, control de versiones de dependencias y
aislamiento de procesos. El archivo \texttt{docker-compose.yml} define los
servicios necesarios:

\begin{itemize}[leftmargin=1.5em]
  \item \textbf{Orquestador}: imagen Python 3.12 con FastAPI y dependencias
  de Spark (PySpark). Expone el puerto 8000 para la API REST de control.
  \item \textbf{Spark Driver}: instancia de Apache Spark en modo local
  (\texttt{local[*]}) que ejecuta los jobs definidos por el orquestador.
  \item \textbf{Red interna}: los contenedores se comunican mediante una red
  Docker bridge que no expone puertos de base de datos al exterior.
\end{itemize}

\subsubsection{Stack Tecnológico del Pipeline}

El procesamiento analítico se organiza en capas que siguen el flujo de datos
desde la extracción hasta el consumo:

\begin{enumerate}[leftmargin=1.5em]
  \item \textbf{FastAPI Orchestrator}: expone una API REST con autenticación
  mediante clave API (\textit{API key}). Recibe solicitudes de ejecución de
  jobs, crea registros de trazabilidad en la tabla
  \texttt{AnalyticsJobRun} y coordina la ejecución de trabajos Spark.
  \item \textbf{Apache Spark (modo local)}: motor de procesamiento batch
  que ejecuta las transformaciones definidas en los jobs. Opera en modo
  \texttt{local[*]} para aprovechar los dos núcleos disponibles, con
  almacenamiento intermedio en disco para datasets que excedan la memoria
  asignada.
  \item \textbf{Supabase PostgreSQL}: base de datos transaccional de la cual
  se extraen los datos operacionales mediante JDBC. Las tablas Gold
  resultantes se escriben de vuelta en el mismo clúster para consulta por
  parte del dashboard administrativo.
\end{enumerate}

\subsubsection{Arquitectura Medallion Aplicada}

El pipeline implementa la arquitectura Medallion descrita en la sección
anterior, con las siguientes responsabilidades por capa:

\begin{itemize}[leftmargin=1.5em]
  \item \textbf{Bronze} (\textit{extracción cruda}): los jobs de extracción
  leen las tablas transaccionales desde PostgreSQL mediante el conector JDBC
  de Spark. Los datos se almacenan en formato Apache Parquet particionados
  por fecha de ejecución, preservando el esquema original sin
  transformaciones.
  \item \textbf{Silver} (\textit{enriquecimiento}): los jobs de
  transformación aplican joins entre dominios, normalizan estados
  (\texttt{PENDING}, \texttt{IN\_PROGRESS}, \texttt{DONE},
  \texttt{CANCELLED}), convierten marcas temporales a zonas horarias
  locales y preparan campos derivados para agregación.
  \item \textbf{Gold} (\textit{métricas agregadas}): los jobs de agregación
  producen tablas consolidadas con indicadores de negocio como ventas por
  franja horaria, productos más solicitados, tiempos de preparación promedio
  y ocupación de mesas. Estas tablas se escriben en PostgreSQL para consumo
  directo por el dashboard.
\end{itemize}

% ============================================================
\subsection{Componentes del Pipeline}
% ============================================================

El pipeline analítico se descompone en componentes autónomos que pueden
ejecutarse de manera independiente o encadenada. Cada componente se
materializa como un \textit{job} de Spark invocado por el orquestador
FastAPI.

\subsubsection{Orquestador FastAPI}

El orquestador constituye el punto de entrada al pipeline. Expone una API
REST con dos endpoints principales:

\begin{itemize}[leftmargin=1.5em]
  \item \texttt{POST /jobs}: recibe una solicitud de ejecución con el tipo
  de job, el identificador del restaurante y los parámetros de la ventana
  de datos. Valida la clave API presente en el encabezado
  \texttt{X-Analytics-Key}, crea un registro en
  \texttt{AnalyticsJobRun} con estado \texttt{PENDING} y lanza el job
  correspondiente mediante \texttt{spark-submit}.
  \item \texttt{GET /jobs/\{id\}}: devuelve el estado actual del job
  identificado, incluyendo el timestamp de inicio, duración estimada,
  progreso de capas Medallion y cualquier error registrado durante la
  ejecución.
\end{itemize}

El orquestador mantiene un registro de cada ejecución en la tabla
\texttt{AnalyticsJobRun}, que incluye los siguientes campos:

\begin{itemize}[leftmargin=1.5em]
  \item \texttt{jobRunId}: identificador único generado como UUID v4.
  \item \texttt{jobType}: tipo de job (\texttt{BRONZE\_EXTRACTION},
  \texttt{SILVER\_ENRICHMENT}, \texttt{GOLD\_AGGREGATION},
  \texttt{DEMAND\_FORECASTING}, \texttt{GOLD\_WRITEBACK},
  \texttt{VELOCITY\_HOURLY}).
  \item \texttt{restaurantId}: identificador del restaurante procesado.
  \item \texttt{status}: estado del job (\texttt{PENDING},
  \texttt{RUNNING}, \texttt{COMPLETED}, \texttt{FAILED}).
  \item \texttt{startedAt}, \texttt{completedAt}: marcas temporales de
  inicio y finalización.
  \item \texttt{errorMessage}: mensaje de error en caso de fallo.
\end{itemize}

\subsubsection{Job 1: Extracción Bronze}

El primer job del pipeline ejecuta la extracción de datos transaccionales
desde Supabase PostgreSQL hacia la capa Bronze. Utiliza el conector JDBC
de Spark para leer tablas completas o incrementales según la ventana de
datos solicitada.

\begin{itemize}[leftmargin=1.5em]
  \item \textbf{Entrada}: tablas transaccionales de Supabase
  (\texttt{Order}, \texttt{OrderItem}, \texttt{TableSession},
  \texttt{Product}, \texttt{Category}, \texttt{Station},
  \texttt{Restaurant}, entre otras).
  \item \textbf{Salida}: archivos Apache Parquet con particionamiento
  por fecha (\texttt{partitionBy("execution\_date")}), almacenados en
  \texttt{/data/bronze/YYYY-MM-DD/}.
  \item \textbf{Configuración}: la conexión JDBC utiliza la URL
  definida en la variable de entorno \texttt{SUPABASE\_JDBC\_URL},
  con credenciales almacenadas en el archivo \texttt{.env} del VPS y
  montadas en el contenedor Docker como variables de entorno.
\end{itemize}

El particionamiento por fecha permite la ejecución de cargas
incrementales, donde únicamente se extraen registros nuevos o
modificados desde la última ejecución, optimizando el uso de
recursos y el tiempo de procesamiento.

\subsubsection{Jobs 2 a 6: Silver, Gold y Procesamiento Avanzado}

Los jobs subsecuentes completan el flujo Medallion y añaden
capacidades de análisis avanzado:

\begin{enumerate}[leftmargin=1.5em]
  \item \textbf{Job 2 --- Enriquecimiento Silver}: lee los datos Bronze,
  aplica joins entre dominios transaccionales, normaliza estados y
  convierte marcas temporales. Produce datasets normalizados en
  \texttt{/data/silver/YYYY-MM-DD/}.

  \item \textbf{Job 3 --- Agregación Gold}: consume los datasets Silver
  y ejecuta agregaciones predefinidas: ventas totales por sucursal y
  franja horaria, productos más solicitados, ticket promedio por mesa,
  tiempos de preparación por estación y rotación de mesas. Los resultados
  se almacenan en \texttt{/data/gold/YYYY-MM-DD/}.

  \item \textbf{Job 4 --- Pronóstico de Demanda}: aplica modelos de
  series temporales sobre las métricas Gold históricas para proyectar la
  demanda esperada por franja horaria y día. Los resultados se almacenan
  en tablas Gold de pronóstico, permitiendo al dashboard mostrar
  estimaciones para planificación operativa.

  \item \textbf{Job 5 --- Escritura de Vuelta (\textit{Write-Back})}:
  toma las tablas Gold generadas en Parquet y las escribe en las tablas
  analíticas de PostgreSQL mediante el conector JDBC. Cada registro
  incluye el \texttt{jobRunId} que lo generó, garantizando trazabilidad
  completa desde la extracción hasta el resultado final.

  \item \textbf{Job 6 --- Velocidad Horaria}: calcula indicadores de
  velocidad operativa por franja horaria: órdenes por hora, tiempo
  promedio de atención, tasa de ocupación de mesas y productividad por
  estación de cocina. Estas métricas permiten al administrador
  identificar cuellos de botella operativos.
\end{enumerate}

\subsubsection{Monitoreo del Pipeline}

La visibilidad y trazabilidad de las ejecuciones del pipeline se
materializan a través de los siguientes mecanismos de monitoreo:

\begin{itemize}[leftmargin=1.5em]
  \item \textbf{Registro en base de datos}: al iniciar y finalizar cada
  job, el orquestador actualiza el registro correspondiente en la tabla
  \texttt{AnalyticsJobRun} con el estado, duración, filas procesadas y
  cualquier error detectado. Esta tabla constituye la fuente canónica
  para consultar el historial de ejecuciones mediante el endpoint
  \texttt{GET /jobs/\{id\}}.
  \item \textbf{Archivos de registro (\textit{logs})}: los contenedores
  Docker del orquestador y de Spark escriben registros detallados en la
  salida estándar, capturados por el driver de logging de Docker. Estos
  logs pueden consultarse mediante \texttt{docker logs} e incluyen
  información de diagnóstico como trazas de excepciones, tiempos de cada
  etapa de Spark y estadísticas de registros leídos y escritos.
  \item \textbf{Métricas de Spark}: el contexto de Spark expone métricas
  de ejecución a través de su interfaz web en el puerto 4040, accesible
  desde el VPS. Estas métricas incluyen el progreso de stages, uso de
  memoria del driver y distribución de tareas entre los núcleos
  disponibles.
\end{itemize}

Este esquema de monitoreo proporciona un registro completo y consultable
de cada ejecución del pipeline, facilitando la detección de fallos
recurrentes, el análisis de tendencias de rendimiento y la generación de
reportes para el administrador del sistema.

\subsubsection{Topología de Archivos del Pipeline}

El sistema de archivos del VPS organiza los artefactos del pipeline bajo
una estructura predecible:

\begin{lstlisting}[basicstyle=\ttfamily\small,breaklines=true]
/srv/bouquet-pipeline/
  docker-compose.yml
  .env
  orchestrator/
    main.py
    requirements.txt
  spark-jobs/
    bronze_extraction.py
    silver_enrichment.py
    gold_aggregation.py
    demand_forecasting.py
    gold_writeback.py
    velocity_hourly.py
/data/
  bronze/YYYY-MM-DD/
  silver/YYYY-MM-DD/
  gold/YYYY-MM-DD/
\end{lstlisting}

% ============================================================
\subsection{Configuración del Entorno en VPS}
% ============================================================

La puesta en marcha del pipeline en el VPS requiere la configuración
de servicios del sistema operativo, variables de entorno y límites de
recursos. A continuación se detallan los componentes de configuración
necesarios.

\subsubsection{Registro de Ejecuciones}

Cada invocación del pipeline queda registrada de forma permanente en la
tabla \texttt{AnalyticsJobRun} de la base de datos PostgreSQL. Esta
tabla funciona como bitácora central de todas las ejecuciones y es
consultada por el dashboard administrativo para desplegar el historial
de jobs, sus estados y métricas de rendimiento. Al residir en la misma
base de datos que los datos transaccionales y analíticos, el registro de
ejecuciones no requiere infraestructura externa adicional y se beneficia
de las políticas de respaldo y replicación ya existentes en Supabase.

\subsubsection{Orquestación de Contenedores con Docker Compose}

El archivo \texttt{docker-compose.yml} define la configuración de los
servicios del pipeline. Los aspectos principales incluyen:

\begin{itemize}[leftmargin=1.5em]
  \item \textbf{Política de reinicio}: configurada como
  \texttt{restart: unless-stopped} para ambos servicios, garantizando
  que los contenedores se reinicien automáticamente tras una caída del
  proceso o un reinicio del servidor.
  \item \textbf{Montaje de volúmenes}: el directorio
  \texttt{/data} del host se monta en ambos contenedores para
  compartir los archivos Parquet generados. El archivo
  \texttt{.env} se monta como volumen de solo lectura para las
  variables de entorno.
  \item \textbf{Red interna}: los servicios comparten una red Docker
  tipo \texttt{bridge} denominada \texttt{pipeline-net}, que permite
  la comunicación entre el orquestador y el driver de Spark sin
  exponer puertos al exterior, salvo el puerto 8000 del orquestador
  para la API REST.
\end{itemize}

\subsubsection{Variables de Entorno}

Las credenciales y parámetros de configuración se almacenan en el
archivo \texttt{.env}, ubicado en el directorio raíz del proyecto en
el VPS. Las principales variables definidas son:

\begin{itemize}[leftmargin=1.5em]
  \item \texttt{SUPABASE\_JDBC\_URL}: URL de conexión JDBC a la base
  de datos Supabase, con formato
  \texttt{jdbc:postgresql://<host>:<port>/<database>}.
  \item \texttt{SUPABASE\_USER} y \texttt{SUPABASE\_PASSWORD}:
  credenciales de solo lectura para la extracción de datos.
  \item \texttt{ANALYTICS\_API\_KEY}: clave secreta utilizada para
  autenticar las solicitudes a la API REST del orquestador.
\end{itemize}

El archivo \texttt{.env} se excluye del control de versiones mediante
\texttt{.gitignore} y se transfiere al VPS por un canal seguro
(SCP/SSH).

\subsubsection{Límites de Recursos para Spark}

Dado que el VPS dispone de 8 GB de RAM y 2 núcleos virtuales, se
configuran los siguientes límites para Apache Spark:

\begin{itemize}[leftmargin=1.5em]
  \item \texttt{SPARK\_DRIVER\_MEMORY=2g}: asigna 2 GB de RAM al
  proceso driver de Spark, responsable de coordinar la ejecución de
  jobs y mantener el contexto de la aplicación.
  \item \texttt{SPARK\_SHUFFLE\_PARTITIONS=4}: limita el número de
  particiones durante las operaciones de \textit{shuffle} (joins,
  agrupaciones) a 4, evitando la creación excesiva de archivos
  temporales y la saturación del sistema de archivos.
  \item \texttt{SPARK\_LOCAL\_DIR=/tmp/spark-temp}: define el
  directorio para archivos temporales de Spark, evitando que ocupen
  el almacenamiento principal.
\end{itemize}

Estos valores se establecen como variables de entorno en el
\texttt{docker-compose.yml} y se aplican automáticamente al iniciar
los contenedores.

% ============================================================
\subsection{Seguridad del Pipeline}
% ============================================================

La seguridad del pipeline analítico se aborda en múltiples capas:
autenticación de la API, aislamiento de red, protección del servidor y
trazabilidad de ejecuciones.

\subsubsection{Autenticación por Clave API}

El orquestador FastAPI implementa un middleware de validación que
inspecciona el encabezado \texttt{X-Analytics-Key} en cada solicitud
entrante. El mecanismo de autenticación funciona de la siguiente manera:

\begin{itemize}[leftmargin=1.5em]
  \item Toda solicitud a los endpoints \texttt{POST /jobs} y
  \texttt{GET /jobs/\{id\}} debe incluir el encabezado
  \texttt{X-Analytics-Key} con un valor que coincida con la variable
  de entorno \texttt{ANALYTICS\_API\_KEY}.
  \item Si el encabezado está ausente o el valor no coincide, el
  orquestador responde con código HTTP 401 (\textit{Unauthorized}) y
  un mensaje descriptivo.
  \item La validación se realiza mediante comparación en tiempo
  constante para mitigar ataques de temporización
  (\textit{timing attacks}).
  \item La clave API se rota periódicamente mediante actualización del
  archivo \texttt{.env} y reinicio del contenedor del orquestador.
\end{itemize}

Este esquema proporciona una capa de seguridad adecuada para un
pipeline interno, donde el único consumidor de la API es el dashboard
administrativo de Bouquet.

\subsubsection{Aislamiento de Conexiones a Base de Datos}

Las conexiones a la base de datos PostgreSQL se rigen por el principio
de mínimo privilegio y aislamiento de red:

\begin{itemize}[leftmargin=1.5em]
  \item \textbf{Conexión JDBC interna}: la URL de conexión JDBC utiliza
  el hostname interno de Supabase, no expuesto públicamente. Spark se
  conecta a través de la red Docker interna del VPS, sin que el puerto
  de base de datos sea accesible desde Internet.
  \item \textbf{Usuario de solo lectura}: la extracción Bronze utiliza
  un usuario de base de datos con permisos exclusivamente de lectura
  (\texttt{SELECT}) sobre las tablas transaccionales. Las escrituras de
  tablas Gold utilizan un usuario separado con permisos de escritura
  limitados al esquema analítico.
  \item \textbf{Sin exposición de puertos}: el archivo
  \texttt{docker-compose.yml} no publica puertos de base de datos hacia
  el host. La comunicación entre el contenedor de Spark y PostgreSQL se
  realiza exclusivamente a través de la red externa del VPS, gestionada
  por el firewall.
\end{itemize}

\subsubsection{Trazabilidad de Ejecuciones}

La tabla \texttt{AnalyticsJobRun} proporciona un registro de auditoría
de todas las ejecuciones del pipeline, permitiendo reconstruir la
trazabilidad completa de cada dato analítico generado:

\begin{itemize}[leftmargin=1.5em]
  \item Cada registro en las tablas Gold incluye el campo
  \texttt{jobRunId} que lo generó, estableciendo una cadena de
  trazabilidad desde el resultado agregado hasta la ejecución
  específica del pipeline.
  \item El esquema de base de datos garantiza la integridad referencial
  entre \texttt{AnalyticsJobRun} y las tablas de resultados mediante
  claves foráneas.
  \item Las marcas temporales \texttt{startedAt} y
  \texttt{completedAt} permiten auditar la ventana de procesamiento de
  cada job y detectar posibles solapamientos entre ejecuciones.
  \item El campo \texttt{errorMessage} preserva el diagnóstico de
  fallos, facilitando la depuración retrospectiva de incidencias.
\end{itemize}

Este mecanismo de trazabilidad, implementado enteramente sobre la capa
de persistencia de PostgreSQL, permite al administrador del sistema
verificar la integridad de los datos analíticos y diagnosticar fallos
sin depender de servicios externos.

\subsubsection{Endurecimiento del VPS}

El servidor VPS se protege mediante las siguientes medidas de seguridad
a nivel de sistema operativo:

\begin{itemize}[leftmargin=1.5em]
  \item \textbf{Acceso SSH}: configurado exclusivamente mediante
  autenticación por clave pública. El archivo
  \texttt{/etc/ssh/sshd\_config} deshabilita el acceso por contraseña
  (\texttt{PasswordAuthentication no}) y el acceso como superusuario
  (\texttt{PermitRootLogin prohibit-password}).
  \item \textbf{fail2ban}: servicio activo que monitorea los registros
  de SSH y bloquea temporalmente direcciones IP que excedan el umbral
  de intentos fallidos de autenticación. La configuración por defecto
  aplica un baneo de 10 minutos tras 5 intentos fallidos en una ventana
  de 600 segundos.
  \item \textbf{Firewall con iptables}: las reglas de firewall
  restringen el tráfico entrante a los puertos estrictamente
  necesarios: 22 (SSH), 80/443 (HTTP/HTTPS para el dashboard) y 8000
  (API del orquestador, con restricción por IP de origen cuando es
  viable).
  \item \textbf{Actualizaciones automáticas}: el paquete
  \texttt{unattended-upgrades} aplica parches de seguridad del sistema
  operativo de forma automática, reduciendo la ventana de exposición a
  vulnerabilidades conocidas.
\end{itemize}

% ============================================================
\subsection{Resultados y Métricas}
% ============================================================

El despliegue del pipeline analítico en el VPS produce resultados
tangibles en dos dimensiones: la generación de datos analíticos para el
dashboard administrativo y el registro de trazabilidad de las
ejecuciones.

\subsubsection{Generación de Datos Analíticos}

Una vez que el pipeline se encuentra operativo, se espera la generación
de los siguientes artefactos de datos:

\begin{itemize}[leftmargin=1.5em]
  \item \textbf{Archivos Parquet diarios}: cada ejecución diaria del
  pipeline produce conjuntos de archivos Parquet en las tres capas
  Medallion, organizados por fecha en la estructura
  \texttt{/data/\{bronze|silver|gold\}/YYYY-MM-DD/}. Estos archivos
  permiten la reconstrucción de cualquier estado analítico histórico
  sin necesidad de reprocesar desde el origen transaccional.
  \item \textbf{Tablas Gold en PostgreSQL}: los resultados de las
  agregaciones se persisten en tablas del esquema analítico de
  Supabase, listas para ser consultadas por el dashboard mediante
  Prisma. Cada registro incluye el identificador del job que lo generó
  (\texttt{jobRunId}), permitiendo correlacionar métricas con
  ejecuciones específicas del pipeline.
  \item \textbf{Métricas clave}: el dashboard administrativo consume
  indicadores como ventas totales por sucursal y franja horaria,
  productos más solicitados, ticket promedio, tiempos de preparación
  por estación, tasa de ocupación de mesas y pronóstico de demanda.
\end{itemize}

\subsubsection{Registro de Ejecuciones}

La tabla \texttt{AnalyticsJobRun} constituye el registro canónico de
todas las ejecuciones del pipeline. Para cada job ejecutado se registra:

\begin{itemize}[leftmargin=1.5em]
  \item Identificador único del job (\texttt{jobRunId}).
  \item Tipo de job ejecutado.
  \item Restaurante procesado.
  \item Estado final (\texttt{COMPLETED} o \texttt{FAILED}).
  \item Marcas temporales de inicio y finalización.
  \item Duración total en segundos.
  \item Mensaje de error en caso de fallo.
\end{itemize}

Este registro, almacenado directamente en PostgreSQL junto con los
datos analíticos, permite al administrador monitorear la salud del
pipeline, detectar degradaciones de rendimiento y auditar la integridad
de los datos generados.

\newpage
